\section{Restrictions on extension kernels}\label{sec:obstructions}

We derive necessary conditions from the kernel and image of the extension
map. Throughout this section, \(S\) is a stabilizer with
\(S\cap V_T=R\), \(S^\perp\cap W=C\), and extension map \(M\).
Every nonzero kernel element represents a logical operator supported on
\(T\).
There are two kernels to distinguish. The kernel \(K=\ker M\) records
logical operators supported on \(T\). The kernel \(N=\ker\tau\)
records stabilizer words supported on the complementary coordinates:
their projected local class is zero, and an element of \(R\) removes
the remaining local representative. The hull identities below relate
these two kinds of local information.

\begin{proposition}[Kernel and hull identities]\label{prop:kernel}
For an extension, put \(K=\ker M\), \(q=\dim K\), and let \(t\) be the
rank of the symplectic form restricted to \(K\). Then
\begin{align}
 \rank\Gamma&=a-2q+t,&
 h&=b-a+2q-t,\label{eq:hull}\\
 N:=\ker\tau&\subseteq\rad(P),&
 \dim N&=b-a+q.\label{eq:N}
\end{align}
Every nonzero \(e\in K\) has \(\cost(e)\geq d(S)\), even without
purity. Moreover \(q-t/2\leq k\).
\end{proposition}
\begin{proof}
The adjoint map has image \(K^\perp_E\), of dimension \(a-q\).
To see this directly, \(e\) annihilates \(\im\tau\) if and only
if \(\langle e,\tau(x)\rangle=0\) for every \(x\in P\), or
equivalently \(Me=0\). Nondegeneracy of \(E\) then gives
\(\im\tau=K^\perp_E\).
The radical of its restricted form is \(K\cap K^\perp\), of dimension
\(q-t\) by rank--nullity on the restricted form. Pullback along the surjection \(P\to K^\perp\) leaves
the form rank unchanged, giving \(a-q-(q-t)\). Subtracting from \(b\)
gives \eqref{eq:hull}. The map \(\tau\) has rank \(a-q\), giving
\eqref{eq:N}; a vector in its kernel pairs trivially with all of \(P\).

For \(e\ne0\) in \(K\), a prefix representative with zero suffix lies
in \(S^\perp\). It cannot lie in \(S\), because \(S\cap V_T=R\).
Its minimum weight is at least the logical distance. The space
\(K\) embeds in the \(2k\)-dimensional nondegenerate logical space
\(S^\perp/S\). A subspace of dimension \(q\) and form rank \(t\)
has an isotropic subspace of dimension \(q-t/2\); that dimension is at
most \(k\).
\end{proof}

\subsection{Counting cosets of small weight}

Let \(Q=\im M\). The syndrome subspace \(Q\) is naturally
\(N^\perp_W/C\): these are precisely the suffix classes annihilating
\(N=\ker\tau\). The map \(M\) induces a bijection \(E/K\to Q\).
For \(\eta\in E/K\), put
\(\kappa(\eta)=\min_{e\in\eta}\cost(e)\).
\(\kappa(\eta)\) is the least weight available among all local
classes that have the same suffix image. If that suffix image has small
minimum weight, every one of its possible local partners must compensate
with a sufficiently large weight. Counting these partners gives the next
inequality.

\begin{proposition}[Coset-counting inequality]\label{prop:capacity}
In the localized setting of Theorem~\ref{thm:extension}, for
\(0\leq j<d\),
\begin{equation}\label{eq:capacity}
 |\{0\ne u\in Q:\delta(u)\leq j\}|
 \leq |\{0\ne\eta\in E/K:\kappa(\eta)\geq d-j\}|.
\end{equation}
If every nonzero class in \(E\) has minimum weight below \(d\), then
\[
 K=0,\quad h=b-a,\quad N=\rad(P),\quad Q=\im\Gamma.
\]
\end{proposition}
\begin{proof}
For a fixed \(\eta\), all its members have the same syndrome image.
Minimizing \eqref{eq:exact} over them gives
\(\kappa(\eta)+\delta(M\eta)\geq d\).
Each nonzero syndrome class of minimum weight at most \(j\) therefore
corresponds to a distinct nonzero coset in \(E/K\) whose minimum weight
is at least \(d-j\). This proves the inequality.
The minimum-weight hypothesis and Proposition~\ref{prop:kernel} imply
\(K=0\). Then \eqref{eq:hull} and \eqref{eq:N} give the remaining
statements, with \(\im\Gamma=(\ker\Gamma)^\perp\) in syndrome
coordinates.
\end{proof}

Violation of \eqref{eq:capacity} excludes an extension without enumerating
the matrices \(M\). The inequality alone is not sufficient for existence.

\subsection{A binary bound for a rank-one local stabilizer}

\begin{proposition}[Kernel projection bound]\label{prop:griesmer}
Suppose \(R=\langle Z^m\rangle\) and \(d>\lfloor m/2\rfloor\).
The physical \(X\)-projection maps \(K\) injectively to an even binary
linear code of length \(m\), dimension \(q\), and distance at least
\(\Delta=2(d-\lfloor m/2\rfloor)\). Consequently,
\begin{equation}\label{eq:griesmer}
 m\geq\sum_{i=0}^{q-1}\left\lceil\frac{\Delta}{2^i}\right\rceil.
\end{equation}
\end{proposition}
\begin{proof}
Adding \(Z^m\) leaves the \(X\)-mask unchanged, so this projection is
well defined on \(E\). Commutation with \(Z^m\) makes its weight
\(w\) even. If \(z\) counts \(Z\)-letters off that mask, then
\[
 \cost(e)=w+\min(z,m-w-z)\leq\lfloor m/2\rfloor+w/2.
\]
A nonzero class represented by a \(Z\)-only word has minimum weight at most \(\lfloor m/2\rfloor\),
so cannot lie in \(K\). This proves injectivity, and
\(\cost(e)\geq d\) gives \(w\geq\Delta\).
The last inequality is the binary Griesmer bound applied to that projected
code. For completeness, choose a minimum-weight binary word of weight
\(w\). Puncturing its support leaves a code of dimension one less and
distance at least \(\lceil w/2\rceil\): the two representatives of a
nontrivial quotient class have weight sum equal to \(w\) plus twice
their common suffix weight, and each has weight at least \(w\).
Only zero and the chosen word can lie wholly in that support. Induction
gives the stated sum.
\end{proof}

The binary Griesmer bound applies to the projected kernel. For \(m=8,d=6\), it gives \(q\leq4\).
Section~\ref{sec:seventeen} sharpens this to three using parity and the
logical symplectic dimension.

\subsection{Balanced local planes and redundancy}

When the local stabilizer is generated by two full-support words, the
minimum-weight condition determines the symbol counts of every nonzero
kernel class. This gives the following restriction.

\begin{theorem}[Balanced-plane obstruction]\label{thm:balanced}
Suppose \(S\cap V_T=R=\langle X^m,Z^m\rangle\), where \(m=4\omega\),
\(\omega\geq1\), and \(d(S)\geq3\omega\). Then \(2^q\leq m\).
If \(\omega\) is even, \(K\) is isotropic. If \(\omega\) is odd, then
\(q\leq2\), with restricted form rank zero for \(q\leq1\) and two
for \(q=2\). In particular,
\[
 \begin{array}{ll}
 n-k\geq2m-2-\lfloor\log_2m\rfloor,&\omega\text{ even},\\
 n-k\geq2m-4,&\omega\text{ odd}.
 \end{array}
\]
No purity assumption is needed.
\end{theorem}
\begin{proof}
For a word on \(T\), let \(n_I,n_X,n_Y,n_Z\) be its four symbol
counts. Addition of \(R\) permutes the four counts, so the minimum
weight modulo \(R\) is \(m-\max n_L\leq3m/4\). A nonzero kernel class has minimum weight
at least \(3\omega\); hence every count is \(\omega\). All four words in that
class have weight \(3\omega\).

Let \(\widetilde{K}\) be the preimage of \(K\) in \(R_T^\perp\), and put
\(L=|\widetilde{K}|=2^{q+2}\). Its weight enumerator is
\(1+(L-4)z^{3m/4}+3z^m\). The second MacWilliams coefficient of its
dual equals
\[
 \frac{6m^2}{L}-\frac{3m}{2}.
\]
This coefficient counts weight-two words in the dual code. It is
nonnegative, so \(L\leq4m\), or \(2^q\leq m\).

The Pauli parity function \(\varepsilon(v)=\wt(v)\bmod2\) satisfies
\begin{equation}\label{eq:polar}
 \varepsilon(u+v)=\varepsilon(u)+\varepsilon(v)+\langle u,v\rangle.
\end{equation}
If \(\omega\) is even, all words of \(\widetilde{K}\) are even, so \(\widetilde{K}\), and hence
\(K\), is isotropic. If \(\omega\) is odd, two distinct nonzero kernel
classes pair to one, since the three representatives involved have odd
weight. Three independent classes would contradict bilinearity:
\(\langle u,v+w\rangle=1+1=0\), although the same rule requires one.
\(q\leq2\), with the asserted ranks.

Here \(a=2m-4\), \(b=n-k-2\). In the even case,
\(\rank M=a-q\leq b\) gives \(n-k\geq2m-2-q\).
In the odd case, the smallest possible \(\rank\Gamma=a-2q+t\)
is \(a-2\), so \(b\geq a-2\), giving the second bound.
\end{proof}

For example, the specified local plane on twelve coordinates at distance
at least nine requires \(n-k\geq20\). The hypothesis fixes the entire stabilizer supported on those twelve
coordinates; the conclusion applies to codes with that local structure.

\subsection{Choosing a tractable support}

\begin{proposition}[Support selection]\label{prop:support}
Let \(B(f,d)\) be an upper bound on the binary dimension of an additive
length-\(f\), minimum-distance-\(d\) code. If the stabilizer of a pure
\(\qc{n,k,d}\) code contains a nonzero word with support \(T\),
where \(|T|=m\), then \(r=\dim(S\cap V_T)\) satisfies
\[
 1\leq r\leq B(n-m,d)-n-k+2m.
\]
\end{proposition}
\begin{proof}
The shortened centralizer has dimension \(n+k-2m+r\) by
Theorem~\ref{thm:extension}'s projection argument, and distance at least
\(d\) by purity. Apply the additive dimension bound.
\end{proof}

The selected word need not have minimum stabilizer weight. This freedom
is useful when an exact enumerator certificate forces a particular weight
whose complement is already classifiable.
